%%% LaTeX template file for UP poster.
%%% ------------------------------------------------------------------------------------------------

\documentclass[pageborders=False, colorscheme1, a0paper, 35pt]{upposter}
% Drawing page borders is useful when the poster will be printed
% on a larger paper size, and needs to be cut afterwards.
% In other cases it is best to omit option pageborders.

% \usepackage[ngerman]{babel} % german (hyphenation... )
\usepackage[english]{babel} % english (hyphenation... )
\usepackage[utf8]{inputenc}
\pdfcompresslevel=0

% TODO
\title{Hardware-Accelerated Load Balancing on Bluefield-3 DPUs - XenoFlow Implementation and Performance Evaluation}
%%% List the authors here.
% TODO
%\author{M. Itarbeiter}
%%% List the cooperations and funding partners here.
% TODO
\institute{\begin{tabular}{ll}
\huge Max Schrötter, Sten Heimbrodt, Bettina Schnor
\end{tabular}}

% Block Style without borders
\defineblockstyle{Plain}{
	titlewidthscale=1,
	bodywidthscale=1
}{
	\setlength{\blocklinewidth}{0pt}
}


\begin{document}
	%%% Uncomment this to add the mathpluslogo
	% TODO 
	
	%%% This command creates your title.
	\maketitle % Create title

	
	
	
	%%% Here you can put your poster content.
	%%%
	%%% A box with a title is created by the following command:
	%%% \block{Title}{Content}
	%%% For a box without a titlebar just leave the title empty:
	%%% \block{}{Content}
	%%%
	%%% The content of a block allows regular latex code (with some exceptions).
	%%%
	%%% Blocks can span the width of the poster or can be placed inside columns.
	%%% Columns are defined like this:
	%%%
	%%% \begin{columns}
	%%% \column{.3} % this is a column spanning 30 percent of the width of the poster.
	%%% \block{Title1}{Content1}
	%%% \block{Title2}{Content2}
	%%% \column{.7} % this is a column spanning 70 percent of the width of the poster.
	%%% \block{Title3}{Content3}
	%%% \end{columns}
	%%%
	%%% At the very end of the poster creating process you should make sure that your
	%%% blocks at the end of the poster align. You can for example play around with
	%%% \vspace{2cm}
	%%%
	%%% Assuming your file is named 'poster.tex' you can compile it by opening a commandline,
	%%% navigating to the folder containing the file and typing
	%%%    pdflatex poster.tex
	%%%
	%%% A word to the papersize:
	%%% The frames at ZIB have b1paper format. Since the fontsizes and measurements should be uniform,
	%%% please use the following doumentclass and don't change neither the colorscheme nor
	%%% the fontsizes or fontfamilies:
	%%% \documentclass[grey,b1paper,fontserif]{zibposter}
	%%%
	%%% If you or your conference prefer a0paper or you like a different colorscheme
	%%% you can change these parameters in the documentclass.
	
	%%% --- begin content ------------------------------------------------------------------------------
%	\begin{columns}
%		
%		\column{0.55}
%		
\useblockstyle{Plain}
\block{}{
	\begin{tikzfigure}
		\includegraphics[width=0.7\linewidth]{graphics/demo/Datacenter_3.png}
	\end{tikzfigure}
}

	
%		
%	\end{columns}
%	

\useblockstyle{upposter}
\block{Motivation}{
	Some short explanation of the goals or findings\\
	
	\begin{minipage}[t]{.3\linewidth}
		\begin{itemize}
			\item Point 1
			\item Point 2
			\item Point 3
			\item ...
			\item Point 41
			\item Point 42
			\item Point 43
		\end{itemize}
	\end{minipage}
	\begin{minipage}[t]{.7\linewidth}
		\begin{itemize}
			\item Goal 1
			\item Goal 2
			\item Goal 3
			\item ...
			\item Goal 41
			\item Goal 42
			\item Goal 43
		\end{itemize}
	\end{minipage}
}

	
	\begin{columns}
		
		\column{.55}
		
		\block{A picture}{
			\begin{tikzfigure}[Ligand Excluded Surfance \cite{ref1}]
				\includegraphics[width=0.7\linewidth]{graphics/demo/ligandExcludedSurface}
			\end{tikzfigure}
		}
		
		\block{Box Title}{
			Some text ...

			\vspace*{160mm}
		}
		
		\column{.45}
		
		\block{Box Title}{
			A box with some latex math. First Euler's equation
			\[ e^{i\pi} + 1 =0 \]
			Second the geometric series
			\[ \sum_{k=0}^{n-1} q^k = \frac{1-q^n}{1-q} \]
		}
		
		\block{}{
			A box without a title.  A box without a title.

			\vspace*{207mm}
		}
		
		\block{Font awesome}{
			Note that you can use fontawesome symbols:
			\begin{itemize}
				\item[\faUniversity] a university
				\item[\faBook] a book
				\item[\faGlobe] a globe
				\item[\faUsers] users
			\end{itemize}
			and much more, for references search the internet
		}
		
	\end{columns}
	
	\begin{columns}
		
		\column{.55}
		
		\block{Publications}{
			\begin{thebibliography}{9}
				\bibitem{ref1} N. Lindow, D. Baum, H.-C. Hege. Ligand Excluded Surface: A New Type of Molecular Surface.
				\textit{IEEE Transactions on Visualization and Computer Graphics,} Vol. 20, No. 12, pp. 2486--2495, 2014.    
			\end{thebibliography}

			\vspace{20mm}
		}
		
		\column{.45}
		
		\block{Funded by}{			
			Our awesome partner and ...
			\hfill
			%\includegraphics[height=2.4cm]{graphics/logos/ZIB-logo}
			\vspace{33mm}
		}
		
	\end{columns}
	
	%%% --- end content --------------------------------------------------------------------------------
\end{document}
